% =====================================================================
%  R. Nielsen, ET PAR ORD I ANLEDNING AF PROF. SCHARLINGS APOLOGIE FOR DR. MARTENSENS DOGMATIK.
%  Kjøbenhavn: C. A. Reitzel (Bianco Luno), 1850.  24 pp.  Fraktur.
%
%  TRANSCRIPTION -- GENERATED; DO NOT EDIT BY HAND.  Rebuild, chained (&&), from the repo root:
%    python3 ebook-method/scharling-apologie/draft.py --emit --force
%      this template + tesseract's script/Fraktur reading of KB's scan (the scan has NO text layer; the reading is
%      set as one by ocr.py MODEL=layer); ~/bibliotek/Nielsen, Rasmus/1850-scharling-apologie.pdf, see SCANS.tsv;
%    python3 ebook-method/collate/check.py scharling-apologie --witness ocr --apply
%      the decisions in decisions.tsv, taken against an independent reading by a different model (tesseract deu_latf);
%    python3 ebook-method/scharling-apologie/oepass.py --apply
%      the ø pass and the word fixes of the unattested-word sweep (word-fixes.tsv);
%    python3 ebook-method/scharling-apologie/finish.py --apply
%      the footnotes, corrections across markup, Latin (\textit), letterspacing (\emph), PRINTED AS IS marks.
%
%  WORD-ACCURACY (2026-10-06).  No e-book, no earlier transcription and no text layer exist.  The draft is
%  tesseract's script/Fraktur reading (tessdata_best), set as a layer; it was collated word by word, with
%  punctuation, quote kinds and word division, against tesseract's deu_latf model: 830 disagreements in 5705
%  words.  About 540 were settled by rule (rules.py, each decision named): each model's known Fraktur faults on
%  this print undone (Fraktur reads k as f, sk as f or st, v as y, x as r; deu_latf has no æ), accepted only
%  where the result is a word attested in at least two other books of the corpus, two attested words going to
%  the image unless one is OCR junk beside a far commoner word; a third reading (tesseract dan) voting only for
%  an attested word.  The rest were read at the image: 278 crops read BLIND by three readers (controls 6 of 6),
%  decided by script (readsbox.py) or by the main session from the readers' transcriptions.
%  Ø: all three models read the print's ø as o.  oepass.py restores it where the o-form is no word (122 tokens)
%  and, for pairs that are both words (for/før, bor/bør, tor/tør, fore/føre ...), from 56 tight crops of the
%  word alone, without context (controls 4 of 4): 10 ø.
%  Then: the unattested-word sweep (94 words found in no other book: 43 OCR faults with one possible reading set
%  by sense, 15 read at the image, the rest genuine or misprints); two page readers on every page (letterspacing, antiqua,
%  paragraph starts, every footnote read whole), and the main session's own look at every page for
%  letterspacing (it found „tænker“, p. 4, that the readers missed) and at 2x zoom for the doubtful words.  The
%  layer dropped no text (check.py --dropped).  Shared faults both models read alike were set from the page
%  images (finish.py GLOBAL and POST: iffe, fun, fal, theologiste, Loftestang, Parador, Dine, Dieblik ...).
%  NOT ESTABLISHED: misprints that both models read alike as real words, beyond those met in the sweep and the
%  page looks (no full letter-by-letter image pass); letterspacing of single short words.
%  PAGE MARKERS: 22 of 22 exact (check.py --pages).
%
%  CONVENTIONS (as Dr. H. Martensens dogmatiske Oplysninger belyste).  Body in Fraktur; aa (not å); ø as printed.
%  Emphasis is letterspacing -> \emph{}.  Latin set in antiqua -> \textit{}.  Greek in Greek type (textalpha).
%  Fraktur I/J are one sort: I for I-words (Idee, Ikke), J for J-words (Jeg).  The double hyphen at a line end is
%  resolved; inside a compound it is "-".  "% --- p. N ---" + \opage{N} at the first word of printed page N.
%  Footnotes *), **) per page, as printed.  Folios, signature marks and the library's stamps and pencil marks are
%  not transcribed.  Print governs; misprints are kept and marked PRINTED AS IS.
% =====================================================================
\documentclass[12pt,a4paper,openany]{book}
\usepackage[utf8]{inputenc}
\usepackage[T1]{fontenc}
\usepackage{libertinus}
\usepackage{libertinust1math}
\usepackage{textalpha}
\usepackage{graphicx}
\DeclareUnicodeCharacter{0254}{\reflectbox{c}}  % ɔ (U+0254) = reversed c
\usepackage[danish]{babel}
\usepackage[protrusion=true,expansion=true]{microtype}
\usepackage{setspace}
\usepackage[margin=1.3in]{geometry}
\usepackage{fancyhdr}
\usepackage{hyperref}
\hypersetup{pdftitle={Et Par Ord i Anledning af Prof. Scharlings Apologie for Dr. Martensens Dogmatik}, pdfauthor={Rasmus Nielsen}, hidelinks}
\setlength{\headheight}{15pt}
\pagestyle{fancy}
\fancyhf{}
\fancyhead[LE,RO]{\thepage}
\fancyhead[C]{\textit{R. Nielsen: Et Par Ord i Anledning af Prof. Scharlings Apologie (1850)}}
\renewcommand{\thefootnote}{\ifcase\value{footnote}\or *)\or **)\or ***)\else
  \arabic{footnote})\fi}
\setstretch{1.3}
\usepackage{marginnote}
\newcommand{\opage}[1]{\marginnote{\footnotesize\textit{[#1]}}}
\newcommand{\apage}[1]{\marginnote{\footnotesize\textit{[#1]}}}

\begin{document}

% --- front matter, read at the image 2026-10-06 (main session).
% p. [1], title.  "Lic.", "Dr." and "1850." in antiqua; "Nielsen," and "Reitzel." letterspaced.  The ornamental rule
% is drawn as a plain rule; KB's stamp and the pencil "Kgl Bibl" are not transcribed.
\begin{titlepage}
\begin{center}
\vspace*{3em}
{\Large Et Par Ord}\\[1em]
i Anledning af\\[1em]
{\Huge Prof. Scharlings Apologie}\\[1em]
for\\[1em]
{\LARGE Dr. Martensens Dogmatik}\\[1em]
af\\[1em]
{\large Lic. R. \emph{Nielsen,}}\\
{\small Prof. i Phil.}\\[2em]
\rule{8em}{0.4pt}\\[2em]
{\large Kjøbenhavn.}\\[0.5em]
Hos Universitetsboghandler C. A. \emph{Reitzel.}\\[0.5em]
{\small Trykt hos Kgl. Hofbogtrykker Bianco Luno.}\\[0.5em]
1850.
\end{center}
\end{titlepage}
% p. [2], the motto, in Greek, at mid-page (the epigram of Demodocus); the attribution in bold antiqua, as printed.
\thispagestyle{empty}
\vspace*{12em}
\begin{center}
\begin{minipage}{0.8\textwidth}
\textit{Μιλήσιοι ἀξύνετοι μέν οὐκ εἰσί· δρῶσι δὲ, οἷαπερ οἱ ἀξύνετοι.}\\[0.4em]
\hspace*{\fill}\textbf{Demodacus.} % PRINTED AS IS: Demodacus (for Demodocus), checked at the image
\end{minipage}
\end{center}
\clearpage


% ===== TEXT (pp. 3--24) =====
% --- p. 3 ---
\setcounter{footnote}{0}%
\opage{3}I „nyt theologisk Tidsskrift“ (1ste Bind, 2det Hefte S.
348—375) er Hs. Høiærværdighed Dr. theol. Prof. C.
E. Scharling --- mærkværdigt nok! --- optraadt som
Apologet for Dr. Martensens Dogmatik.

At et saa energisk Skridt fra Prof. Scharlings
Side er kommet mig noget uventet og tillige har bragt mig
i stor Forlegenhed, vil jeg ikke fordølge. Paa et
skarpt videnskabeligt Indlæg fra Dr. Martensen selv har
jeg længe været forberedt, men paa en Apologie fra Prof.
Scharling, nei, det tilstaaer jeg, paa en saa gemytlig
Apologie var jeg aldeles ikke forberedt.

Hvorfor jeg kalder den Scharlingske Apologie \emph{gemytlig},
er let at forklare. Et videnskabeligt Forsvar
holder sig strengt til Modstanderens Indsigelser og Grunde;
men slige formelle Hensyn ere naturligviis af underordnet
Betydning i en Apologie, hvis hele Beviiskraft lader
sig henføre til Kategorierne: Velvillie og Uvillie, Bifald
og Mishag, Roes og Dadel.

Hvor Prof. Scharling for Tiden har sit videnskabelige
Standpunkt, og om han overhovedet har noget, skal jeg
ikke kunne afgjøre; derimod er det af det foreliggende
% --- p. 4 ---
\setcounter{footnote}{0}%
\opage{4}Aktstykke ingenlunde vanskeligt at angive hans apologetiske
Standpunkt: Prof. Scharling er for nærværende Tid ---
gemytlig rosende Prof. Martensen, gemytlig dadlende
hans Modstandere. Hvad Prof. Scharling nu egentlig
\emph{tænker} om „Troens“ Forhold til „Viden“, om Religionens
Forhold til Philosophien, Christendommens til Theologien,
vil den Smule Skarpsindighed, der blev mig forundt,
ikke kunne udfinde af hans gemytlige Apologie;
hvad jeg derimod uden Anstrengelse kan udfinde, og hvad
der unegtelig ogsaa er lykkedes den theologiske Professor
at fremstille med sjelden theologisk Klarhed, er det for
„den theologiske Videnskab“ vistnok heller ingenlunde uvigtige
Vendepunkt, at Hs. Høiærværdighed er bleven vred
paa Magister Stilling og mig.

Saa bliver da det første facultetsmæssige Udbytte af
den første facultetsmæssige Apologie for vor nyeste Litteraturs
første facultetsmæssige Dogmatik at udtrykke og
fastholde i den ligesaa vægtige som vigtige Sætning: \emph{Professor Scharling er bleven vred paa Magister}
\emph{Stilling og Professor Nielsen}.

At jeg nu ved denne alvorlige og høitidelige Leilighed
først nævner Mag. Stillings Navn, har sin Grund
i Beskedenhed, og skeer alene for Høfligheds Skyld; thi
ifølge Vredens ulige, skjøndt dog vistnok billige Fordeling,
tilkommer det egentlig mig at indtage den første Plads.
Sagen er denne: Paa Magister Stilling er Professor
Scharling kun vred af een Aarsag, den nemlig, at Magister
Stilling har kunnet tillade sig at skrive imod Dr.
Martensen; paa mig er Hs. Høiærværdighed derimod vred
af to Aarsager, thi jeg har ikke alene været saa dristig
at skrive imod Dr. Martensen, men desuden i Forveien
% --- p. 5 ---
\setcounter{footnote}{0}%
\opage{5}været saa dumdristig at skrive Adskilligt imod „de theologiske
Middelmaadigheder.“\footnote[1]{Jfr. „Evangelietroen og den moderne Bevidsthed“.}

Der er da i den gemytlige Apologie en fælles Revselse,
der overgaaer os begge. Vi have begge svarligen
forsyndet os imod den theologiske Videnskab ved „den ubehagelige
Overraskelse“, vore „udførligere Bedømmelser“ af
„den christelige Dogmatik“ have beredet Dr. Martensen;
„den haanende Ironie, den drillende Spot, hvormed Forkastelsesdommen
er ledsaget, den grusomme Lyst, hvormed
vi“ --- at ogsaa jeg skal udtrykke mig gemytlig --- „afskyde Ironiens Pile“ imod den stakkels „berømte“ Forfatter,
„synes (S. 351) ikke at kunne udelukke Formodningen
om, at andre Motiver end den paaberaabte rene
Iver for Sandheden have gjort sig gjældende ved vore
Kritikers Affattelse“; kort sagt: vi blive begge revsede og
det tilgavns, thi vi have begge været slemme.

Men foruden den fælles theologiske Revselse, der
saaledes tilkjendes os begge, har Prof. Scharlings apologetiske Gemytlighed endnu en Revselse, et Overskud, en
Tilgift af Revselse, som særlig tilkjendes mig. Uagtet vi,
som sagt, begge ere slemme, er Magister Stilling dog
alligevel ikke nær saa slem som jeg. Rigtignok er Mag.
Stilling voldsom --- „maaskee endog den voldsomste“ ---
men desuagtet (jfr. Anm. S. 360) „bærer dog Magister
Stillings Skrift ved Siden af personlig Ømfindtlighed
Præget af en subjectiv Overbeviisning, der vækker Læseren
til Alvor og Sympathie“\dots{} Anderledes er det med
mig! „Anderledes er det med Prof. Nielsen. Hans
Fremstilling fremkalder det Indtryk, som om han nok
% --- p. 6 ---
\setcounter{footnote}{0}%
\opage{6}kunde gjøre Tilbudet: „Vil I nu tage min Mening, saa
vil jeg igjen forsvare Jer Mening“\dots{} --- Ak, hvad skal
jeg svare? Bedre var det maaskee, om jeg slet ikke svarede,
men taug stille, og bøiede mig i al Ærbødighed
for den theologiske Revselse! thi, alvorlig talt, hvad skal
jeg svare? Med mig gaaer der til det Yderste; til min
Forbauselse og høist ubehagelige Overraskelse erfarer jeg
nu, at jeg staaer ved Randen af tvende Yderligheder,
tvende forfærdelige Yderligheder, hvis Navne jeg end ikke tør
nævne, men maa overlade Læseren at udfinde af Prof.
Scharlings Beskrivelse.

Den første Yderlighed beskrives (S. 364 Anm.) saaledes:
„Han (Prof. N.) nævner ikke gjerne Theologerne
uden med et eller andet haanende Tillægsord (Skoletheologer,
Middelveistheologer \&c.) % the Fraktur r-rotunda abbreviation ꝛc. (= et cetera), read 2c. by tesseract; rendered \&c. as in propaedeutiske-logik; alt hvad Dumt og Forkeert,
Theologer kunne have forskyldt og ikke forskyldt, samler
han til Træk i de Carrikaturbilleder, hvormed han søger
at gjøre sine theologiske Modstandere til Latterlighedens
Martyrer. Dette hænger forresten sammen med en forunderlig
Illusion om den Anseelse, som han mener, at
Theologerne, især theologiske Doctorer og Professorer, nyde.
Derfor forfølger han disse med den samme Bitterhed og
Satire, som Voltaire (sic!) i sin Tid Præsterne.“

Den anden, endnu langt betænkeligere Yderlighed
beskrives saaledes: Efterat Magister Stillings Sætning:
\textit{quamvis absurdum sit, credo}, foreløbig er berørt, hedder
det i Texten (S. 364): „ikke desto mindre maa der dog
til Grund for den sidste Erklæring, naar den ikke uden
videre skal afvises som udelukkende al videre fornuftig
Overveielse, ligge den Forestilling, at de christelige Dogmer,
hvor urimelige de end fremstille sig for Subjectet,
ikke desto mindre have en objectiv gyldig og evig Sand%
% --- p. 7 ---
\setcounter{footnote}{0}%
\opage{7}hed, hvilken det kun ikke er muligt paa Grund af Forhindringer
eller Indskrænkninger hos Subjectet“ (---\textit{NB}---)
„i en given Tilstand at erkjende.“ Paa denne psychologiske
Forklaring følger nu Noten, og Noten maa man
endelig ikke overspringe, thi Noten er det Bedste, thi
Noten er endnu mere gemytlig end Texten, thi i
Noten ligger Pointen, thi i Noten siger Professor
Scharling: „Om dette (— nemlig hvad der anføres
om Mag. Stillings Sætning —) imidlertid skulde
være Prof. Nielsens Mening, veed jeg ikke; thi han finder
Stillings anførte Formel endnu for tvetydig, og vil have
den udtydet i Sætningen: \textit{credo, quia absurdum est.}
Saaledes kunde her være indtraadt i Virkeligheden, hvad
J. P. Lange (Christl. Dogmatik, 1849, 1ster Th. S.
409) fornylig endnu har sat som \emph{Mulighed}, „daß am
Ende auch noch ein Philosoph kommen könnte, der das
Denken oder das Bewustseyn selbst ebenfalls für eine Illusion
erklären wollte. Was wäre nun da anzufangen?
Man mußte ihn mit dem Zugeständniß, daß sein eignes
Bewustseyn in Taumel gerathen sein könne, nach Hause
schicken.“

Nach Hause! altsaa: nach Hause! ja, lad det saa
blive derved: nach Hause!

Da jeg med min „Evangelietro“ og min „undersøgende
Anmeldelse“ rykkede frem mod en i Litteraturen
vistnok meget anseelig, men dog efter min Overbeviisning
aldeles uholdbar theologisk Videnskab, var jeg endda ikke
forberedt paa saa faa Muligheder; men, at denne \emph{Mulighed},
denne hersens \emph{Mulighed}, hvorom Prof. Scharling
her taler, just skulde være saa nær ved at „være indtraadt
i Virkelighed“, nei, det tilstaaer jeg, derpaa var jeg
% --- p. 8 ---
\setcounter{footnote}{0}%
\opage{8}slet ikke forberedt: det er da ogsaa et „overraskende“ Krigspuds
af den gemytlige Apologie.

Ved dette Krigspuds maa jeg med mine Skrifter
da vel omtrent være tilmode, som Krigsraadet i en maadelig
Fæstning, der uforvarende overrumples af en stærk
og listig Fiende; man forudseer klart, at Modstand er
forgjæves, og dog, da det jo dog alligevel ikke er en aaben
By, men i al „Middelmaadighed“ en Fæstning, føler man, at
man, idetmindste for Ærens Skyld, maa gjøre Noget. Saaledes,
ja netop saaledes gaaer det nu mig med denne
apologetiske Overrumpling. Prof. Scharling kaster Handsken
for den troende Speculation: jeg føler, at jeg maa
gjøre Noget, og forudseer alligevel, at det ikke kan blive
til Noget.

At jeg nu virkelig trods al min „Pirrelighed“ maa have
en vis Grund til at være forsigtig med Prof. Scharlings
litterære Kjendelser, skal Hs. Høiærværdighed selv kunne bevidne,
naar han tænker sig om. For nogle Aar tilbage, det
var just i det Tidspunkt, da --- som Hs. H. selv siger
S. 352 --- „den saakaldte speculative Theologie her tillands
nød en Anseelse, hvis Glands allerede var begyndt
at blegne i dens oprindelige Hjem,“ udgav jeg en, ikke i
streng hegelsk, men i speculativ-theologisk Aand affattet
Commentar over Pauli Brev til Romerne. Denne Commentar
havde strax det Uheld at tildrage sig Prof. Scharlings
Opmærksomhed, hvorfor han da ogsaa efter sin gemytlige
Sædvane fandt sig beføiet til at afgive en litterær
Kjendelse, for med det Samme at vise mig lidt
tilrette.

Havde jeg nu dengang følt Kald til at indlade mig
med Prof. Scharling, havde Leiligheden været saa gunstig,
som muligt. „I de Dage“ --- siger Hs. H. selv S.
% --- p. 9 ---
\setcounter{footnote}{0}%
\opage{9}353 --- „gik der overhovedet en haard Dom over alle Videnskabsmænd,
som ikke vilde eller ikke kunde sætte sig i
umiddelbar Relation til det gjældende System“; altsaa:
i de Dage, da Raisonneurernes Indsigelse imod den speculative
Erkjendelse blev „afviist med streng Tilretteviisning
for aandelig Indskrænkethed og Ubekjendtskab med
Philosophiens Væsen“, i de Dage var det just, at Prof.
Scharling affærdigede mit Forsøg over Romerbrevet\footnote[1]{Tidsskrift for Litteratur og Kritik, udgivet af F. C. Petersen. 1842. 1ste Bind. S. 210.},
og i de Dage havde jeg da ogsaa, ifølge min særdeles
store „Pirrelighed“, havt den skjønneste Leilighed til at
være --- lidt slem imod Prof. Scharling.

Hvorfor forsømte jeg da, siden min „Pirrelighed“ nu
engang er saa overvæties, den gunstige Leilighed, og
bøiede mig i beskeden Taushed under den gemytlige Revselse? --- Kjære Læser, hvad skal jeg svare? Svarer jeg,
at jeg gjorde det, fordi jeg følte en vis Undseelse ved at
angribe en Mand, som kort tilforn havde været min Lærer:
troer man mig vel ikke; thi Hans Høiærværdighed
sigtede mig jo allerede dengang for ironiserende Anmasselse.
Svarer jeg, at jeg gjorde det, fordi jeg helst vilde
skaane en Mand, som aabenbart ved en reen Feiltagelse
var kommen til at afgive sin Kjendelse i en Sag, der
angik den theologiske Speculation: saa troer man mig vel
heller ikke; thi jeg er jo nu engang mistænkt for „andre
Motiver“. Nu, saa vil jeg da svare, hvad sandt er, og
give et Svar, som Prof. Scharling vist bedre vil troe,
end noget andet; jeg vil da svare ham og sige det reent
ud: jeg gjorde det af Frygt.

% --- p. 10 ---
\setcounter{footnote}{0}%
\opage{10}Hvad der indgjød mig denne Frygt, var ingenlunde
Skarpheden i de Scharlingske Indsigelser (thi de manglede
al Skarphed), heller ikke Vægten af de Scharlingske Grunde
(thi de manglede al Vægt), men en fatal Omstændighed,
den nemlig, at Prof. Scharling i mit ungdommelige Forsøg
overhovedet slet ikke kunde finde --- „Aanden“. Ja,
denne Omstændighed var det just, som allerede dengang
lammede min Kraft, og indjog mig en Skræk, jeg siden
den Tid ikke har forvundet.

Til en grundigere Vurdering af min store Forlegenhed
kan jeg dog maaskee allerbedst føre Læseren ad en
lille Omvei.

Før man overhovedet indlader sig i en litterær Strid,
bør man fornuftigviis see sig for, hvem det er, man faaer
med at gjøre, og, saavidt muligt, søge at sondere de polemiske
Kræfter. Hvad navnlig den videnskabelige Polemik
angaaer, forekommer det mig, som om Forfatterne bedst
kunde inddeles i følgende tre Classer: 1) de Forfattere,
som i Genialitetens og i Frihedens Form ere sikkre paa
Princip og Idee; 2) de Forfattere, der erkjende Principet
i Form af Grundsætning, og derfor maae bøie sig
for den regelrette Conseqvents; 3) de Forfattere, der ikke
gaae af Veien for noget Princip, eftersom de i deres litterære
Gemytlighed egentlig ikke vide, hvad et Princip er.

At nu Forfatterne af første Rang ere farlige Polemikere,
især naar de tillige ere begavede med en saa vittig
Phantasie, at de endog kunne skjule det Alvorliges
Tyngde i Spøgens Lethed, forstaaer sig af sig selv. Jeg
nævner her Mag. Kierkegaard og Prof. Heiberg, thi ved
at nævne disse Navne fritager jeg Læseren og mig selv
fra en prosaisk Beskrivelse af den høiere Polemik.

Til Polemikere af anden Classe henhøre alle dygtige
% --- p. 11 ---
\setcounter{footnote}{0}%
\opage{11}og paalideliae Videnskabsmænd, der umiddelbart søge deres % PRINTED AS IS: paalideliae (for paalidelige; 2x zoom)
Styrke i Beviser og Grunde; disse Forfattere ere i
Regelen bundne til den formelle Alvors prosaiske Form,
og kjendes derfor ogsaa bedst under Navn af „de alvorlige
Forfattere“.

Til Polemikerne af Nr. 3 henhøre alle principfrie,
eller rettere: principløse, men dog alvorlige, ufortrødne og
lærde Skribenter, alle værdige Medlemmer af den værdige
Familie \emph{Wagner}\footnote[1]{Jfr. Gocthes Faust.}, der paa en saa værdig Maade % PRINTED AS IS: Gocthes (for Goethes; 2x zoom)
fremmer Videnskaben i alle Retninger og pryder Litteraturen
i alle Lande. Naar en Forfatter af dette Slægtskab
gjør Mine til at indlade sig i en litterær Strid, især i en
dialektisk Principstrid, maa man være forsigtig, meget forsigtig,
thi dette Slags Folk er yderst farligt. --- Med Polemikerne
af første Rang slipper man endda nogenledes
fra det, naar man forresten har en god Sag og en sund
Tanke (selv om man taber, er Tabet altid Gevinst); men
vee Den, der skal værge sig med et Princip imod dem,
der ikke rammes af noget Princip; vee Den, der skal forsvare
en Tankes Eenhed imod dem, der tage alskens Tanker
og alle Slags Ideer til Indtægt, og just derved paa
den skjønneste Maade bevise, „at Menneskers aandelige
Væsen udgjør en Eenhed“; ja vee Den, der falder i Hænderne
paa en Polemiker af Nr. 3, thi --- med Familien
Wagner er der intet Udkomme.

Læseren aner vist allerede min Forlegenhed. --- Under hvilken Rubrik Prof. S. vil henføre mig, er let at see
af den omineuse Yttring: „nach Hause,“ mig er det derimod
ingenlunde let at see, under hvilken Rubrik jeg i denne
% --- p. 12 ---
\setcounter{footnote}{0}%
\opage{12}Sag skal henføre Professor Scharling. At sætte ham
iblandt Polemikerne af første Rang, gaaer dog vel neppe
an (nu det forstaaer sig, hvis en Anden vil gjøre det,
skal jeg ikke sige Noget derimod, men --- jeg tør ikke).
At henføre Hs. Høiærværdighed under Nr. 3 gaaer endnu
mindre an; det vilde jo være en blodig Fornærmelse imod
en Mand, der ikke blot er Medlem af det theologiske Facultet,
men, hvis jeg ellers husker ret, endog Medlem af
Videnskabernes Selskab. Min alvorlige Bestræbelse maa
altsaa gaae ud paa, at faae ham sat i Classe med de
„alvorlige Videnskabsmænd“ under Rubriken Nr. 2; men
hvor utaknemmelig denne Bestræbelse er, og hvor lidet
Haab jeg denne Gang kan have om at føre den til noget
gunstigt Resultat, vil den Sagkyndige let kunne see
ved en flygtig Gjennemlæsning af hans gemytlige Apologie.

Som et af de Hovedtræk, der især charakterisere
denne Apologie, og hvorved Prof. Scharling med samt
hans theologiske Tidsskrift især synes at befinde sig i --- „en
stadig Fremadskriden“, maa jeg her fremhæve den litterære
Ugeneerthed, et interessant Hovedtræk, der ligeledes hænger
sammen med det Gemytlige.

„Det er besynderligt“, siger Prof. Scharling, idet han
(S. 355) omtaler Mag. Kierkegaard, „at endnu Ingen
har gjort denne Forfatters Skrifter til Gjenstand for en
udførligere Analyse, som de baade i sig selv fortjene, og
hvortil der ved det Bifald, som de have vundet hos Læsere
af høist forskjellige Dannelsestrin, synes at være særlig
Opfordring.“ Da jeg nu heri er aldeles enig med
Prof. Scharling, forekommer det mig netop af den Grund
høist „besynderligt“, at Hs. H. med saa megen Ugeneerthed
har kunnet indlade sig paa en saa temmelig affærdigende
% --- p. 13 ---
\setcounter{footnote}{0}%
\opage{13}Bedømmelse af disse Skrifter, uden at der dog i denne
Bedømmelse er mindste Spor af, at Hs. Høiærvædighed % PRINTED AS IS: Høiærvædighed (for Høiærværdighed)
selv har værdiget dem endog blot en flygtig, endsige da,
en udførligere Analyse.

At Prof. Scharling aldeles ikke har værdiget Mag.
Kierkegaards Bøger nogen Analyse, kan strax sees af de
allerførste Linier, hvori han begynder at recensere deres
Form. Der er vel neppe nogen opvakt Student, som ikke
allerede veed, at de Kierkegaardske Skrifter, hvad Formen
angaaer, dele sig i to Hovedclasser, og at ethvert Skrift
i enhver af disse Classer igjen har en i sin Art uendelig
individualiseret, indtil de fineste Enkeltheder omhyggelig
articuleret og nüanceret Form. Denne Formforskjellighed
generer imidlertid ikke Prof. Scharling. Uden at
bekymre sig om nogle ellers usædvanlige Transformationer
slaaer Hs. Høiærv. uden Omstændigheder alle Bøgerne i
een Bunke, og vedbliver saa i gemytlig Ugeneerthed at
tale om „den formløse Brede, den baade til rette og urette
Tid spøgende Ironie, den forcerede Jagen efter Paradoxer
og Antitheser, og den om ogsaa vittige, dog ofte trættende
Snaksomhed, der synes at forhindre Forfatteren, trods
hans forbausende Productivitet, fra at bringe sine Betragtninger
til nogen tilfredsstillende Afslutning.“ Med
denne afsluttende Opfattelse af den Kierkegaardske Form
(en Opfattelse, man vel allerede før har hørt af „travle
Folk“ i Byen) hænger det da ogsaa uden Tvivl sammen,
at Mag. Kierkegaards Dialektik --- besynderligt nok! ---
hos Prof. Scharling (S. 358) „efterlader en Følelse,
som om hans Fremstilling ikke ret var Alvor, men kun
et Forsøg i den høiere Aandsgymnastik.“\footnote[1]{En „udførligere Analyse“ af en af Pseudonymerne (t. Ex. af Johannes Climacus) vilde, hvis den blot nogenledes kunde lykkes, ganske vist afgive stærke Bidrag til Besvarelsen af det Spørgsmaal: \emph{Hvad er Alvor}?} % PRINTED AS IS: t. Ex. (for f. Ex.), as in Dr. H. Martensens dogmatiske Oplysninger belyste

% --- p. 14 ---
\setcounter{footnote}{0}%
\opage{14}Dog, om Formen, og navnlig om Inderlighedens indirecte
Meddelelsesform i dobbeltreflecterede Kategorier, er
jeg, der selv saa haardt støder an imod Formen, uværdig
til at gaae i Rette med Prof. Scharling, hvis smagfulde
Tidsskrift jo ikke blot smager paa Formen, men ogsaa smager
af Formen; en udførligere Formforhandling vilde her desuden
være paa urette Sted, og føre os for langt bort fra den
apologetiske Gemytlighed. Derimod vilde jeg --- om jeg saa
ogsaa ved denne Leilighed skulde støde lidt an imod Formen
--- dog inderlig gjerne give Prof. Scharling det velmeente
Raad, at han aldrig mere giver sig af med at
vurdere eller kritisere de Ting, han endnu ikke ret har
forstaaet, og til den Ende gjøre ham opmærksom paa nogle
Punkter, hvori hans litterære Ugeneerthed synes at være
af den Natur, at den ikke engang kan undskyldes med
det Gemytlige.

Idet Prof. Scharling. (S. 368) vurderer Magister
Kierkegaards Polemik mod „Systemet“, udtaler han sig
paa en saadan Maade, at Den, der ellers ikke er nærmere
underrettet, nødvendigviis maa antage, at Magister Kierkegaard
efterhaanden er skeiet ud fra det berettigede Standpunkt,
hvormed han begyndte, og er endt i en aabenbar
Eensidighed. Det hedder saaledes: „Saalænge Magister
Kierkegaards Polemik indskrænkede sig til en Bestridelse
af, at Systemet var istand til at træde istedetfor Troen\dots{} --- da var den i sin Berettigelse, og man fulgte den med
Bifald og Interesse. Men altsom den gik ud paa at
fjerne Tro og Viden fra hinanden, desto mere førtes den
% the note *) of p. 13 runs over onto the foot of p. 14 („af Johannes Climacus) vilde ... Hvad er Alvor?“): set in full at its mark
% --- p. 15 ---
\setcounter{footnote}{0}%
\opage{15}over til samme Eensidighed, som den saa kraftig bestred\dots{} Derfor kunde ogsaa den fortsatte Polemik, trods
mangt et aandrigt Tankelyn, ikke rigtigen træffe sin Gjenstand,
jo mere den forlod de mere abstracte Regioner, og
gik ind paa de givne historiske Forhold.“

See saaledes bliver nu Mag. Kierkegaards „forbausende
Productivitet“ med Eet affærdiget af Prof. Scharling.
Hvilket Indtryk maa nu en saadan Bedømmelse
gjøre? --- Paa en aldeles naiv og godtroende Læser maa
Bedømmelsen vel gjøre det Indtryk, at Mag. Kierkegaard
her er stedt for en høiere Domstol, at den theologisk-videnskabelige
Overlegenhed, der alt længe har havt et
vaagent Øie med hans mangehaande „dialektiske Vendinger“,
nu endelig, da det synes at gaae for vidt, omsider
bryder Tausheden og viser det urolige „Genie“ tilbage
inden de rette Grændser. Paa mig gjør Bedømmelsen
imidlertid et andet Indtryk, thi paa mig gjør den det
Indtryk, som om det Hele maaskee kunde være løs Tale.
Det Troesparadox, der fremsættes af „Johannes Climacus“,
er jo væsentlig det selvsamme Paradox, som det, der aander
i „\emph{Frygt og Bæven}“. At Stemningen, Formen og
Udviklingen maa være forskjellig, er en Selvfølge, da nemlig
Johannes \textit{de silentio} expectorerer sig i „\emph{dialektisk}
\emph{Lyrik}“, medens Johannes Climacus derimod leverer et
„\emph{existentielt Indlæg}“.

Jeg udtaler denne Paastand uden Omsvøb, ikke blot
fordi jeg mener at være vis i min Sag, men navnlig
fordi det jo egentlig er Prof. Scharling, hvem man,
hvad Benegtelsen angaaer, med Rette kan paalægge \textit{onus}
\textit{probandi}. Vil altsaa Prof. Scharling staae ved sit Ord,
og kan han klart bevise mig, at Mag. K. har i de
nævnte Skrifter forandret Standpunkt. at det Paradox, % PRINTED AS IS: Standpunkt. at (a point, or a comma without its tail)
% --- p. 16 ---
\setcounter{footnote}{0}%
\opage{16}der stilles af Johannes \textit{de silentio}, er et ganske andet og
i et ganske andet Forhold. til \textit{Absurdum}, end det, der % PRINTED AS IS: the point after Forhold
stilles af Johannes Climacus\footnote[1]{Thi at Johannes \textit{de silentio} taler om Abrahams Tro, Johannes Climacus derimod om den christelige Tro, gjør her, hvor det netop gjælder om at udvise Troens Sphære, ingen Forskjel.}; da, ja da skal jeg ikke
blot være villig i al Gemytlighed at lade mig „nach
Hause schicken“, men endog ovenikjøbet gjøre „Apologetiken“
Afbigt og forgylde „den kritiske Indledning“.

Da det imidlertid let kunde træffe sig, at Prof. S.
blev Beviset skyldig, og just derved kom til at bevise lidt
meer, end han egentlig vilde, skal jeg for dette Tilfældes
Skyld ikke her undlade at anføre Noget, som idetmindste
i mine Øine betydelig taler til hans Undskyldning. Idet
Hs. H. gjennemlæste de „udførligere Bedømmelser“ af
Dr. Martensens Dogmatik, er det uden Tvivl hændet
ham, at han saa vel paa „Dogmatikens“ som i det Hele
paa “Tyeologiens“ Vegne har følt sig yderst ubehagelig % PRINTED AS IS: Tyeologiens (for Theologiens), between two high quotes; 2x zoom
afficeret. I sin bekymrede Stemning er han da ved en
ganske naturlig Idee-Association kommen til at tænke
først paa Mag. Kierkegaards Paradox, saa paa mit Paradox,
saa paa Mag. Stillings Paradox, derefter ved en
ligeledes ganske naturlig Idee-Association paa „Proudhon“,
„Max Stirmer“, ja paa Proudhons og Max % PRINTED AS IS: Stirmer (for Stirner; 2x zoom)
Stimers „Consorter“ tilligemed de øvrige „forfængelige % PRINTED AS IS: Stimers (for Stirners; 2x zoom)
Middelmaadigheder“ (\textit{les mediocrités ambitieuses}), som
„forøge Forvirringen ved at kaste hver sit Paradox ud i
Verden“; det har da i et saa bekymret Øieblik været
ham, som om han saae Pandoras aabne Æske og grandgivelig
saae, at Æsken var fuld af --- lutter Paradoxer;
% --- p. 17 ---
\setcounter{footnote}{0}%
\opage{17}det har da i et saa bekymret Øieblik heller ikke været
Hs. H. let at hitte Rede i Paradoxerne, men alle Paradoxer,
christelige, uchristelige, socialistiske, communistiske,
danske, franske, tydske, spanske, kort: alle Verdens virkelige
og uvirkelige Paradoxer have viist sig for ham i en
saadan Mangfoldighed og forvildet sig saaledes imellem
hverandre, at det saa tilsidst slet ikke har været ham muligt
at adskille, hvilket der var hvilket.

Et Hovedtræk, der ligeledes charakteriserer den Scharlingske
Apologie, er en vis litterær Glemsomhed, en gemytlig
Glemsomhed, der i adskillige Punkter forraader, at
den høiærværdige Forfatter vistnok ikke altid maa have
den troeste Hukommelse. Hvorfor Prof. Scharling ikke
just er nogen synderlig Velynder af den philosophiske
Speculation, bekjender han selv, og jeg behøver altsaa
ikke nærmere at forklare ham Grunden; men, at Prof.
Scharling (hvis han ellers ikke i al Stilhed har forandret
sin tidligere Mening) heller ikke er nogen særdeles
Velynder af den theologiske Speculation, vilde jeg endogsaa
temmelig let kunne bevise ham, hvis jeg blot turde
stole en Smule paa hans litterære Hukommelse.

Den Anskuelse af Forholdet imellem Tro og Viden,
der ligger til Grund for Dr. Martensens Dogmatik, er
ganske vist den samme Anskuelse, som den, der ligger til
Grund for min Udvikling af Pauli Brev til Romerne\footnote[1]{Hermed vil jeg dog ingenlunde have paastaaet, at min Udviklingsform skulde være ligesaa classisk som Dr. Martensens (jeg er nu engang ingen „Genius“), heller ikke har jeg i hiint Ungdomsforsøg kunnet sammenarbeide slet saa mange ueensartede Elementer, som Dr. Martensen med saa megen plastisk Kraft har gjort i det fuldendte Værk, hvori „en mystisk Livsaande udøver en lønlig fængslende Magt over den uhildede Læser“. Hvad jeg paastaaer, er kun dette, at „Grundanskuelsen“ er væsentlig den samme, og vedbliver med denne min Paastand, indtil jeg seer et lille Beviis for det Modsatte.}.
% --- p. 18 ---
\setcounter{footnote}{0}%
\opage{18}At Prof. Scharling imidlertid maa have glemt denne
lille Omstændighed, er klart nok af de Ord, hvormed han
(S 371) imødegaaer Kritikens Indvendinger imod den
store Sætning, at Dogmatikeren er „\emph{paa een Gang sin}
\emph{Videnskabs og sin Kirkes Organ.}“ Det hedder
saaledes: „Men enhver saadan Betragtning, der i den
historiske Udvikling anerkjender noget Relativt og Bevægeligt,
bliver afviist af en Kritik, der, medens den gjør
Andre Bebreidelsen for at være blevet staaende paa det
eengang indtagne Standpunkt, selv ikke har løsrevet sig
fra sit forrige Standpunkt i den Henseende, at den jo fra
dette har beholdt ligesaavel den abstracte Opfattelse af
„absolut Viden“, „objectiv Vished“, „sammenhængende Erkjendelse“, som Forestillingen om den christelige Tro
enten som et afsluttet Indbegreb af visse uforanderlige
Læresætninger eller som en aldeles uklar og ubestemmelig
Følelse“. --- Dersom nu Prof. Scharling, da han nedskrev
disse mærkelige Linier, bare havde husket sig lidt om,
og besindet sig paa det Begreb af „absolut Viden“,
„objectiv Vished“, „sammenhængende Erkjendelse“, jeg f. Ex.
har udviklet i min propædeutiske Logik, \textit{item} paa det theologiske
Erkjendelsesprincip, der gaaer igjennem mit Romerbrev,
vilde H. H. ganske vist have strøget den hele
Passus, som aldeles ikke passende paa mig, skjøndt det
dog netop er mig, der har angrebet hiin store Sætning
af Dr. Martensen. Skulde det nemlig være sandt, hvad
% the note *) of p. 17 runs over onto the foot of p. 18 (from „hvori „en mystisk Livsaande“): set in full at its mark
% --- p. 19 ---
\setcounter{footnote}{0}%
\opage{19}Prof. Scharling paastaaer, at jeg nu har en aldeles abstract
Opfattelse af „absolut Viden“, objectiv Vished“, % PRINTED AS IS: no opening quote before objectiv
„sammenhængende Erkjendelse“, saa maa det dog i det
mindste være usandt, hvad dog Prof. Scharling i samme
Aandedræt ligeledes paastaaer, at jeg i denne Henseende
ikke har kunnet løsrive mig fra mit forrige Standpunkt,
thi i de to anførte Skrifter er min Opfattelse af „absolut
Viden“ og „objectiv Vished“ endda ikke saa ganske
abstract; men det er nu eengang en Fatalitet med den
litterære Glemsomhed!

Dengang Prof. Scharling recenserede min Udvikling
af Romerbrevet, var Aanden i den theologiske Speculation
ham endnu en fremmed og ubekjendt Størrelse,
og derfor kunde han --- til min Forbauselse ---
da heller ikke finde Aanden. Nu derimod, da Dr.
Martensen „med sædvanlig Klarhed har udtalet sig i Indledningen
til sin Dogmatik paa en Maade, som i det
Hele vil findes ligesaa tiltalende som tilfredsstillende for
den, der ikke har opgivet Overbeviisningen om, at Menneskets
aandelige Væsen udgjør en Eenhed“, nu faaer
Prof. Scharling pludselig Syn paa Aanden, anerkjender
Principet, og optræder saa i al Gemytlighed som Apologet
for den theologiske Speculation.

Men hvorledes skal jeg nu forklare mig dette Særsyn?
Jeg antager ingenlunde, at Prof. Scharling forsætlig
har villet være partisk eller gjøre mig Uret; meget
mere troer jeg at have opdaget et Synspunkt, hvorfra
det Hele lader sig forklare paa en ligesaa naturlig som
gemytlig Maade. Pointen ligger i, hvad H. H. kalder
„de aandelige Bevægelsers Constellation“. Saalidet Prof.
Scharling nogensinde kan antages at have været partisk
imod „Speculationen“, saa lidet kan han endnu antages
% --- p. 20 ---
\setcounter{footnote}{0}%
\opage{20}at være partisk imod „Paradoxet“. Prof. Scharling holder
vist i Grunden ligesaa meget af „Paradoxet“ som af
„Speculationen“, og ønsker kun, at de kjønt skulle skjønne
derpaa og forliges smukt indbyrdes. Naar Prof. Scharling
altsaa for Tiden taler lidt haardt til „Paradoxet“
og udtrykker sig lidt anderledes om „Speculationen“, end
han forhen har gjort, da er Aarsagen hertil neppe at
søge i nogen væsentlig Forandring i Prof. Scharlings
„Grundanskuelse“; men „her tilbyder sig strax en naturligere
og værdigere Forklaring, naar vi gaae tilbage i Tiden,
og agte paa Forskjellen i de aandelige Bevægelsers
Constellation.“

I de Dage, da Troesspeculationen var saa „hovmodig“,
og saae saa fornemt ned paa den „kritiske Indledning“,
i de Dage, da der overhovedet gik „en haard
Dom over alle Videnskabsmænd og Skribenter, som ikke
vilde eller ikke kunde“ --- give sig af med Speculationen,
i de Dage var Prof. Scharling kun lidet fornøiet med
det hele Speculationsvæsen, enten det saa for Resten
var philosophisk eller theologisk. Ikke uden en vis indre
Tilfredsstillelse hilsede Hs. H. saa de Dage, da en ligesaa
berettiget, som uimodstaaelig Polemik reiste sig imod
det besværlige System, og gjorde „Speculationen“ komisk;
forsaavidt fulgte Prof. S. da „Paradoxet“ med „Bifald og
Interesse“ --- Men saa kom igjen de Dage, da „Speculationen
krøb til Korset, og næsten grædefærdig over
„Paradoxets Vilkaarlighed“ udøste sin Kummer i „den
kritiske Indlednings“ følelsesfulde Barm; og de Dage ere
vel saa omtrent disse Dage, og i disse Dage har Prof.
Scharling da ogsaa følt, at det endelig gaaer forvidt, og
desaarsag besluttet at sige et opmuntrende Ord til den
% --- p. 21 ---
\setcounter{footnote}{0}%
\opage{21}dybt nedbøiede „Speculation“ og holde Styr paa det
uartige „Paradox.“

Ingenlunde skal man derfor beskylde Prof. Scharling
for den Anmasselse at ville udjævne Differenserne i de modstridende
Anskuelser eller, som det hedder, „mediere“ Speculationen
med Paradoxet; nei, saa høit sigter Prof. S.
ikke; han tænker kun som saa: enhver af Potenserne kan
være god for sig, naar de kun i gjensidig Anerkjendelse
ikke genere hinanden eller hindre hinandens „Berømmelse“.
Til Prof. Scharlings Berømmelse være det da sagt, at
man, ifølge hans gemytlige Apologie, ingenlunde bør gaae
strengt i Rette med de modsatte Opfattelser og de ueensartede
Principer, men derimod (S. 373) „bidrage man
til Anerkjendelse og Paaskjønnelse af enhver Bestræbelse,
som med Begeistring for sit Maal og med beskeden Erkjendelse
af de Kræfter, der ere forundte til dets Opnaaelse,
gaaer ud paa at stille de Sandheder, hvorpaa
de enkelte Menneskers og den hele Menneskeslægts evige og
timelige Vel beroer, i et klarere Lys, at begrunde dem
dybere, og paa denne Maade bidrage til at vække og
bevare Ærefrygt og Hengivenhed for deres velgjørende
Indflydelse.“

At der nu ogsaa virkelig kan ligge en smuk Følelse
til Grund for en saa gemytlig Formaning, maa indrømmes;
det er barnlige Forestillinger; men Vanskeligheden er, at
slige barnlige Forestillinger Intet bevise i en videnskabelig
Discussion.

Dersom jeg imidlertid paastod, at der i den hele gemytlige
Apologie intetsteds var gjort nogen Anstalt til at
opklare det videnskabelige Stridspunkt, vilde jeg dog gaae
forvidt. Der er virkelig et Sted, hvor Hs. H. kommer
i Tanker om, at han dog ogsaa burde angive Pro%
% --- p. 22 ---
\setcounter{footnote}{0}%
\opage{22}blemet; dette Sted findes S. 367, og lyder saaledes:
„Det er klart, at Løftestangen i den Dialektik, som her
benyttes, er Spændingen af Begrebet: „\emph{Viden,}“ „\emph{Vished,}“
„\emph{sammenhængende Erkjendelse}“, til en saadan
Grad af „Absoluthed“ og „Objectivitet“, at den enten
ingen Anvendelse finder“ o. s. v.

Altsaa i den urimelig overdrevne Spænding af Begrebet:
„Viden“, „Vished“, „sammenhængende Erkjendelse“,
vil Prof. Scharling see den „Løftestang“, hvormed
en altfor nærgaaende Dialektik har forsøgt at løfte
„den christelige Dogmatik“ af de speculative Hængsler.
Hvorledes det nu er gaaet til, at Prof. Scharling just har
faaet Syn paa dette Begreb, og saa troet, at det var Løftestangen,
er ikke vanskeligt at gjætte. Da Prof. Martensen
rykkede frem med sin store Dogmatik, og vilde gjøre
Plads til Høire og Venstre, var Begrebet: „sammenhængende
Erkjendelse“, netop det Varsko, hvormed han betydede „de absolute i Troen“ og „de absolute i Viden“,
at de nu maatte gaae af Veien, og ikke komme ham for
nær med deres „Strøtanker og Aphorismer, Indfald
og Glimt.“ Da jeg nu i denne Trængsel blev betagen
af Forskrækkelse, og var i Angst for at blive ynkelig
overkjørt, greb jeg efter „Løftestangen“, og fandt den heldigviis
i „Paradoxet“ hos „Johannes Climacus“. For
imidlertid at minde Læseren om den forblommede Udfordring,
der syntes at lyde saa kjækt i det Martensenske Slagord,
lod jeg i min „undersøgende Anmeldelse“ Begrebet:
„sammenhængende Erkjendelse“, løbe med som en lille
glaneur, at der i Tilfælde af, at „Løftestangen“ gjorde
sin Virkning, ikke skulde spildes noget Ax af det speculative
Læs. Hvorledes jeg for Resten indbilder mig at
have brugt „Løftestangen“, er tydeligt nok af min sidste
Thesis: „\emph{den christelige Dogmatik er princip}%
% --- p. 23 ---
\setcounter{footnote}{0}%
\opage{23}\emph{løs}“. --- At nu Prof. Scharling ikke er falden paa at
søge „Løftestangen“ i denne Thesis, forekommer vistnok høist
besynderligt, da Stangen jo ellers er følelig nok; men, som
sagt, jeg frygter meget, at Hs. H. har sine egne Meninger
om, hvad et Princip er.

At jeg paa ingen Maade har kunnet bruge Begrebet:
„sammenhængende Erkjendelse“, som „Løftestang“,
vil Enhver indsee, der overhovedet veed, hvad en Løftestang er.
En Løftestang maa jo være fast og stærk; men Begrebet: „sammenhængende
Erkjendelse“, er just et meget blødt og bøieligt
Begreb, der snoer sig i alle mulige Grader. Der
gives jo mange Slags Sammenhæng og mange Grader
i enhver Sammenhæng; der gives reel Sammenhæng,
ideel Sammenhæng, fast Sammenhæng, løs Sammenhæng:
Alt, selv den gemytlige Apologie har sin Sammenhæng; hvis Prof. Scharling altsaa i Kategorien:
Sammenhæng, mener at have fundet en Stang, kan det
visselig ikke være en „Løftestang“; det maa da vel være --- en anden Stang; dog vel aldrig en Liimstang!

Til Løn for sin ædle Stræben at kuske „Paradoxet“
og holde Orden i Litteraturen, venter Prof. Scharling,
at „de Mænd, mod hvis Adfærd han nu har erklæret
sig“, skulle kjølne „deres Harme i en Philippika ---
hvorved han forudseer en ikke mindre ublid Behandling,
end den, Marsyas erfarede af den fortørnede Gud.“ ---
Ak! hvor det dog i denne Verden gaaer underlig omkring
med den menneskelige Forudseen. Jeg forudsaae nu saavist,
at jeg skulde faae et Svar fra Prof. Martensen,
og see, saa faaer jeg --- en Apologie fra Prof. Scharling;
Prof. Scharling forudseer saa igjen, at han vil
faae „en Philippika“ fra mig, og see, saa faaer han
--- et Par gemytlige Ord.

Det er da for Resten besynderligt nok, at Prof. S. har
kunnet „forudsee“, at hans gemytlige Apologie virkelig skulde
% --- p. 24 ---
\setcounter{footnote}{0}%
\opage{24}fortjene en Philippika, og derved tillige forudsee, at det vilde
gaae ham som Marsyas; denne Forudseen maa da vist ogsaa
hidrøre fra en af de mørke Muligheder, han stundom
seer, naar han seer Paradoxerne i Pandoras Æske.
Hvad mig anbelanger, kan Hs. H. saamænd være rolig
for, at jeg aldrig skal hjemsøge ham med nogen Philippika,
eller kjølne min Harme paa en uskyldig Mand, der
vist aldrig vil vække min Harme. Jeg har i mine Skrifter
aldrig yttret nogen Harme mod Prof. Scharling, aldrig
udfordret Prof. Scharling, og, reentud sagt, sjeldent % PRINTED AS IS: reentud (for reent ud)
tænkt paa Prof. Scharling. Vil Prof. S. desuagtet
godhedsfuld vedblive, som hidtil, at tænke paa mig, vise
mig lidt til Rette og tage sig af min Opdragelse, skal
jeg ogsaa vide at paaskjønne det i al Stilhed; kun maa
Prof. S. ikke tage det ilde op, om han kun sjeldent seer
skriftligt Beviis paa min Taknemmelighed; jeg har, som
sagt, en Frygt for at indlade mig formeget med Prof.
Scharling, der altid ængster mig med sine mørke Muligheder.


At nu Mulighederne kunne være mørke nok, skulde
jeg, der ikke er sikker fra nogen Side, men kun subsisterer
ved at troe det Bedste og frygte det Værste, da ogsaa
allermindst fordriste mig til at benegte. Men selv
om det Værste skulde skee, selv om det engang skulde
ske, at enten „Paradoxet“ eller „Speculationen“ fandt
sig beføiet til at paategne mig mit Skudsmaal, give
mig Afsked som en ubrugelig Tjener, og stikke mig ---
„nach Hause“; saa vilde jeg dog, idet jeg stikkes „nach
Hause“, ret gjerne medtage den Trøst, at „den theologiske
Videnskab“ omsider maatte have fundet sit Princip, og
derhos tillige bevaret den Overbeviisning, at Professor
Scharlings „aandelige Væsen udgjør en Eenhed.“

\begin{center}\rule{4em}{0.4pt}\end{center} % a small ornamental rule closes the text

\end{document}
